\documentclass[11pt,a4paper]{article}
\usepackage[margin=2.2cm]{geometry}
\usepackage{fontspec}
\setmainfont{TeX Gyre Termes}\setsansfont{TeX Gyre Heros}\setmonofont{Latin Modern Mono}[Scale=0.9]
\usepackage{booktabs,longtable,array,enumitem,xcolor,hyperref,titlesec,fancyhdr,microtype}
\hypersetup{pdftitle={The Open Software Stack Across All Repositories of Naziru Halilu},pdfauthor={Naziru Halilu},colorlinks=true,linkcolor=blue!50!black,urlcolor=blue!50!black}
\titleformat{\section}{\Large\bfseries\sffamily}{S\thesection}{0.6em}{}
\titleformat{\subsection}{\large\bfseries\sffamily}{S\thesection.\arabic{subsection}}{0.5em}{}
\setlength{\parskip}{4pt}\setlength{\parindent}{0pt}
\pagestyle{fancy}\fancyhf{}\fancyfoot[C]{\thepage}\renewcommand{\headrulewidth}{0pt}
\newcolumntype{R}[1]{>{\raggedright\arraybackslash}p{#1}}
\newcommand{\gh}[1]{\href{https://github.com/halilunaziru73-creator/#1}{\texttt{\detokenize{#1}}}}
\newcommand{\ghn}[2]{\href{https://github.com/halilunaziru73-creator/#1}{#2}}
\begin{document}
\begin{center}{\large\bfseries\sffamily Software Stack Supplement\par}\vspace{4pt}
{\LARGE\bfseries\sffamily The Open Software Stack Across All Repositories of Naziru Halilu\par}\vspace{6pt}
{\normalsize Languages and libraries and techniques and tools found by reading the code of every repository\par}\vspace{6pt}
{\small Naziru Halilu \quad halilunaziru73@gmail.com \quad \href{https://github.com/halilunaziru73-creator}{github.com/halilunaziru73-creator}\par}\end{center}
\tableofcontents
\clearpage

\section{Purpose and how this was made}
This document is the companion of Supplementary Material 5 of the SWAFL series, widened from one repository to the whole account. Every public repository of the account (17 in all: SWAFL and 14 other research repositories, the personal website, and the profile repository) was copied, and its files, its dependency lists (\texttt{requirements.txt}, \texttt{pyproject.toml}, \texttt{environment} files, PyInstaller \texttt{.spec} files, GitHub workflow files) and its code were read. The Python packages come from the import lines of the code (scripts and notebooks), not from memory; the techniques come from the methods stated in each README and from the names of the functions and classes in the code. The software of the SWAFL repository itself is described in full in Supplementary Material 5 of that repository and is only summarised here.

\textbf{Words used.} A \emph{language} is the way instructions are written. A \emph{library} or \emph{package} is code written by others that the code calls for one job; a package is a library wrapped so that one command installs it (\texttt{pip install}). A \emph{program} or \emph{tool} is a finished application that you start. A \emph{technique} is a method of analysis. Counts are of files and lines on the day of writing; a package is listed as used only when the code imports it.

\section{The repositories at a glance}
\begin{longtable}{@{}R{3.8cm}R{3.3cm}R{3.6cm}R{4.2cm}@{}}
\caption{The 17 repositories, what each is, and its main languages and libraries.}\label{tab:repos}\\
\toprule Repository & What it is & Languages (lines) & Main libraries and tools\\\midrule\endfirsthead
\toprule Repository & What it is & Languages (lines) & Main libraries and tools\\\midrule\endhead
\ghn{Coupled-Water-Fluxes-Pore-Scale-Dynamics-Precision-Irrigation-Mulching-Feedback-Loops}{Coupled Water Fluxes (SWAFL)} & SWAFL Desktop 1.1.0 and the four articles & JavaScript (14\,100) and Python (about 10\,000 over the folders) and \LaTeX\ (103 files) and Markdown & Electron and Leaflet and GeoTIFF.js and SheetJS and numpy and matplotlib and pandas and scipy and scikit-learn and statsmodels and Plotly; Daisy 7.1.14 (C++)\\
\ghn{halilunaziru73-creator.github.io}{Personal website} & the personal website and the CVs & HTML and CSS (3\,800) & Google Fonts; PDF files of the CVs; GitHub Pages\\
\ghn{halilunaziru73-creator}{Profile page} & the profile page & Markdown and HTML & GitHub profile README\\
\ghn{Operationalizing-GIS-and-Machine-Learning-across-Contrasting-Cropping-Systems}{Operationalizing GIS and Machine Learning} & three field campaigns in one validation framework & Python (920) & numpy and pandas and scipy and scikit-learn and matplotlib and Pillow\\
\ghn{Zaria-Crop-ET-and-Irrigation-DSS}{Zaria Crop ET and Irrigation DSS} & crop evapotranspiration and irrigation decision support & Python (7\,200) and JSON & Tkinter and numpy and pandas and matplotlib and python-docx and Pillow and PyInstaller\\
\ghn{NaCROP}{NaCROP} & the Nigeria version of the same crop and irrigation program & Python (9\,600) and JSON & Tkinter and numpy and pandas and matplotlib and python-docx and Pillow and PyInstaller\\
\ghn{Digital-Twin-for-the-Evaluation-of-Experimental-Gully-Biocontrol-Using-Morning-Glory-Ipomoea-spp}{Gully Biocontrol Digital Twin} & a four-layer digital twin of a gully & Python (3\,300) & numpy and pandas and scipy and scikit-learn and matplotlib and joblib and Pillow; setuptools; GitHub Actions\\
\ghn{Geometry-Agnostic-Contrastive-Learning-GACL}{GACL} & a contrastive learning framework for agronomic images & Python (4\,600) and a notebook & PyTorch and torchvision and scikit-image and scikit-learn and pandas and openpyxl and PyQt6 and Pillow and matplotlib and scipy\\
\ghn{Hybrid-CNN-BiLSTM-Attention-Sediment-Transport-Agricultural-Gully-System}{Hybrid CNN BiLSTM Attention} & a hybrid deep network for sediment transport & Python (1\,560) & numpy and pandas and scikit-learn and matplotlib (the network itself is written with numpy)\\
\ghn{Real-Time-RGB-Proxy-Vegetation-Indexing-and-Texture-Analysis-for-UAV-and-Handheld-Crop-Imagery}{RGB Proxy Vegetation Indexing} & a desktop program for crop image analysis & Python (6\,800) & PyQt6 and PyTorch and torchvision and scikit-image and scikit-learn and numpy and pandas and openpyxl and Pillow and matplotlib and rasterio; PyInstaller\\
\ghn{GIS-based_delineation_of_areas_suitable_for_livestock_slurry_application}{Livestock Slurry Suitability} & suitability mapping for slurry application & Python (1\,130) and QGIS project files & rasterio and geopandas and shapely and pyproj and fiona and PyKrige and scikit-learn and TensorFlow with Keras and pandas and numpy and scipy and matplotlib; QGIS\\
\ghn{Integrated-Decision-Support-Model-for-Dairy-Processing-under-Milk-and-Market-Variability}{Dairy Processing DSS} & a decision-support model for a dairy plant & Python (1\,640) & pandas and numpy and scikit-learn and scipy and matplotlib and pytest and flake8; GitHub Actions\\
\ghn{Multitask-Singularity-Regression-Kriging-for-Spatio-Temporal-Soil-NPK-Crop-Stress-Yield-Prediction}{MSRK} & a GeoAI framework for soil and yield maps & Python (1\,640) & scikit-learn and PyKrige and scipy and numpy and pandas and matplotlib\\
\ghn{Geospatial-data-analysis}{Geospatial Data Analysis} & a six-module course in GIS and remote sensing & Jupyter notebooks (2\,600 lines of code) and some JavaScript & rasterio and geopandas and shapely and rasterstats and folium and OWSLib and GDAL (\texttt{osgeo}) and scikit-learn and matplotlib and pandas and numpy; Google Earth Engine\\
\ghn{GIS-Monitoring-of-Vineyard-Ripening-Brix}{Vineyard Ripening Brix} & NDVI and sugar monitoring of a vineyard & QGIS project and style files & QGIS 4 (project and style and metadata files); Sentinel-2 NDVI; spreadsheets\\
\ghn{GIS-Monitoring-of-Grape-Moth-Vineyard-Pest}{Grape Moth Pest} & trap monitoring of a vine pest & QGIS project and style files & QGIS 4; digital elevation model; shapefiles and GeoPackage\\
\ghn{STADOS-Spectral-Temporal-Adaptive-Dark-Object-Subtraction}{STADOS} & atmospheric correction of Sentinel-2 NDVI & Python (410) and a notebook & numpy and pandas and scipy and matplotlib; Jupyter\\
\bottomrule\end{longtable}

\section{Languages and formats}
\begin{longtable}{@{}R{3.4cm}R{4.6cm}R{7.0cm}@{}}
\caption{All languages and formats found in the 17 repositories, with their role.}\label{tab:lang}\\
\toprule Language or format & Where & Role and why it is used\\\midrule\endfirsthead
\toprule Language or format & Where & Role and why it is used\\\midrule\endhead
Python & 11 repositories besides SWAFL (38\,800 lines of Python and 2\,960 lines of notebook code in them) & the language of the analysis and the models and the desktop programs. Why: the libraries for data and maps and learning are mature, and the scripts run unchanged on any computer.\\
Jupyter notebooks (\texttt{.ipynb}) & the course repository and STADOS and GACL & step by step analysis that shows code and results together. Why: a reader can run and follow each step.\\
JavaScript & SWAFL and the course repository & the program SWAFL and the interactive pages. Why: it runs in a browser and in Electron without a compiler.\\
HTML and CSS & the website and the graphical abstracts and the SWAFL window & the pages that people see. Why: opens in any browser with no installation.\\
\LaTeX\ and BibTeX & the SWAFL repository & articles and supplements with exact equations.\\
Markdown & every repository (\texttt{README.md}) & the explanation of each repository. Why: plain text that GitHub shows as a page.\\
JSON & most repositories & configuration and data and results that both people and programs read.\\
YAML (\texttt{.yml}) & workflow files & the instructions of GitHub Actions.\\
TOML (\texttt{pyproject.toml}) & the digital twin and the dairy model & how a Python package is described and built. Why: a package can then be installed by name.\\
CSV and XLSX & most research repositories & tables of data and results. Why: any spreadsheet program opens them.\\
GIS formats: shapefile (\texttt{.shp} with \texttt{.shx} \texttt{.dbf} \texttt{.prj} \texttt{.cpg}) and GeoPackage and GeoTIFF (\texttt{.tif}) and KML and QGIS project (\texttt{.qgz}) and QGIS style (\texttt{.qml}) and QGIS metadata (\texttt{.qmd}) & the GIS repositories & maps and rasters and styles. Why: open standards that QGIS and Python read.\\
Word documents (\texttt{.docx}) & manuscripts and reports & the text of manuscripts for co-authors.\\
Archives (\texttt{.7z}) and model files (\texttt{.joblib}) & the course and the digital twin & large data and trained models.\\
PyInstaller \texttt{.spec} & NaCROP and the crop DSS and the image program & the recipe that turns a Python program into one Windows file.\\
Batch files (\texttt{.bat}) & NaCROP and the crop DSS & start files for Windows.\\
\bottomrule\end{longtable}

\section{Python packages used in all repositories}
\begin{longtable}{@{}R{2.9cm}R{3.6cm}R{8.5cm}@{}}
\caption{Every third-party Python package imported by the code, with its role.}\label{tab:py}\\
\toprule Package (import name) & Repositories & Role and why it is used\\\midrule\endfirsthead
\toprule Package (import name) & Repositories & Role and why it is used\\\midrule\endhead
numpy & 12 repositories & arrays and the numerical core. Why: every other package builds on it.\\
pandas & 12 & tables of observations and results. Why: grouping and joining by name avoids index errors.\\
matplotlib & 12 & every figure and map layout. Why: one library gives print quality figures.\\
scipy & 8 & statistics and curve fitting and image shifting and interpolation.\\
scikit-learn (\texttt{sklearn}) & 9 & random forest and gradient boosting and ridge and elastic net and PCA and k-nearest neighbours and k-means and cross validation. Why: standard and repeatable learners.\\
PyTorch (\texttt{torch}) and torchvision & GACL and the image program & deep networks (hypergraph transformer and ResNet101 and Faster R-CNN). Why: flexible training of new architectures.\\
TensorFlow and Keras & slurry suitability & a multilayer perceptron and a patch-based convolutional network. Why: fast prototyping of small networks.\\
scikit-image (\texttt{skimage}) & GACL and the image program & texture features (grey-level co-occurrence matrix) and image filters.\\
PyKrige (\texttt{pykrige}) & slurry suitability and MSRK & kriging and regression kriging. Why: geostatistical interpolation without writing the variogram code.\\
rasterio & 6 (in some it is optional) & reading and writing rasters (GeoTIFF). Why: the standard Python door to GDAL.\\
GeoPandas and Shapely and pyproj and Fiona & slurry suitability and the course & vector data (points and lines and polygons) and map projections.\\
rasterstats & the course & zonal statistics of rasters inside polygons.\\
folium and OWSLib and animatplot & the course & web maps and web map services and animated plots.\\
GDAL (\texttt{osgeo}) & the course & raster and vector conversion.\\
Pillow (\texttt{PIL}) & 6 & reading and saving images and the icons of the programs.\\
openpyxl & GACL and the image program (listed in the requirements and used by pandas) & lets pandas read Excel files.\\
joblib & the digital twin and GACL & saving and loading trained models.\\
python-docx (\texttt{docx}) & NaCROP and the crop DSS & writing the Word farm reports.\\
xlrd & NaCROP and the crop DSS & reading old Excel files (optional).\\
PyQt6 & GACL and the image program & the GIS-style window of the desktop program.\\
Tkinter (standard library) & NaCROP and the crop DSS and the image program & the simple desktop window. Why: it ships with Python.\\
pytest and flake8 & the dairy model & unit tests and style checks.\\
PyInstaller & NaCROP and the crop DSS and the image program & builds the standalone Windows program. Why: the user then needs no Python.\\
setuptools & the digital twin and the dairy model & packages the code so it installs by name.\\
IPython and Jupyter & the course and STADOS & running the notebooks.\\
\bottomrule\end{longtable}
The standard library is also used (for example \texttt{dataclasses} in 23 files, \texttt{argparse} in 12, and \texttt{threading}). Modules written by the author inside a repository (for example \texttt{bgdt}, \texttt{gacl}, \texttt{nacrop}, \texttt{zaria\_maize}, \texttt{dairy\_dss}, \texttt{msrk\_engine}, \texttt{qgis\_ui}, \texttt{core}, \texttt{cartohelpers}, \texttt{hybrid\_model}) are not third-party packages and are not listed.

\section{Techniques}
\begin{longtable}{@{}R{3.6cm}R{5.0cm}R{6.6cm}@{}}
\caption{The techniques used, where, and why.}\label{tab:tech}\\
\toprule Technique & Where it is used & Role and why\\\midrule\endfirsthead
\toprule Technique & Where it is used & Role and why\\\midrule\endhead
Process simulation with Daisy & SWAFL (10\,190 runs) & the soil and plant and atmosphere model whose runs the suite reads.\\
Evapotranspiration equations (FAO-56 Penman-Monteith and ASCE and Priestley-Taylor and Makkink and Turc and Hargreaves-Samani and Thornthwaite and Blaney-Criddle and Dalton) & the crop DSS and NaCROP & ten ways to estimate crop water use, compared on the same weather.\\
Soil water balance and irrigation scheduling (TAW and RAW depletion) & the crop DSS and NaCROP and SWAFL & when and how much to irrigate. \\
Crop growth in four steps (canopy cover and transpiration and biomass and yield) & NaCROP & yield from water use.\\
Pedotransfer and hydraulic functions (Saxton and Rawls and van Genuchten and Mualem) & SWAFL & soil hydraulics from texture.\\
Hydraulic design (Hazen and Williams and Manning and volume balance) & SWAFL & pipe and surface irrigation design.\\
Bayesian data assimilation and hydrological and erosion modelling (HEC-HMS and HEC-RAS 2D and RUSLE-LISEM named in the README) & the gully digital twin & combines sensors and models into a twin that can be asked what-if questions.\\
Digital twin architecture (physical and digital and brain and service layers) & the gully digital twin & a structure that joins measurement and prediction.\\
Neural network written from scratch (Conv1D and bidirectional LSTM and attention pooling and dense layers with a gradient check) & the sediment transport repository & the layers are coded with numpy so that every step can be inspected.\\
Model explanation (permutation importance and kernel SHAP and Monte Carlo dropout and bootstrap intervals) & the sediment transport and the digital twin and the dairy model & shows which inputs matter and how sure the result is.\\
Contrastive and hypergraph transformer learning & GACL & learns image features that do not depend on camera geometry.\\
Classical image analysis (RGB proxy vegetation indices and NDVI and NDRE and grey-level co-occurrence texture and PCA and k-nearest neighbours and k-means) & the image program and GACL & plant condition from ordinary photographs.\\
Pretrained deep detectors (ResNet101 and Faster R-CNN) & the image program & optional detection in images.\\
Regression kriging with random forest trend & slurry suitability and MSRK and the validation framework & a data-driven trend plus geostatistical interpolation of the residuals.\\
Local singularity and multifractal indices & MSRK & a multiscale feature for each location.\\
Multicriteria overlay (GIS-MCDA) & slurry suitability & the conventional suitability map that the new methods are compared with.\\
Deep surrogates (multilayer perceptron and patch-based CNN) & slurry suitability & continuous suitability maps.\\
Dark-object subtraction and temporal harmonisation & STADOS & corrects the atmosphere in multi-temporal Sentinel-2 NDVI.\\
Optical and radar remote sensing (Sentinel-1 and Sentinel-2) and supervised and unsupervised classification & the course & land cover maps.\\
Terrain analysis (altitude and slope and aspect from a digital elevation model) & grape moth and slurry & shows where conditions differ.\\
Trap monitoring and mating-disruption costing & grape moth & pest pressure by place and the cost of control.\\
Logarithmic regression forecast of sugar (Brix) against day of year & vineyard ripening & predicts the harvest date.\\
Machine learning model comparison (ridge and elastic net and random forest and gradient boosting) with leave-one-out validation & the dairy model & picks a profit model that suits few data.\\
Linear programming and Monte Carlo simulation & the dairy model & allocates capacity among products and shows risk.\\
Tiered validation (reproducibility and mapping and association and prediction and decision and transfer) & the validation framework & states how far each result can be trusted.\\
\bottomrule\end{longtable}

\section{Tools}
\begin{longtable}{@{}R{3.8cm}R{4.8cm}R{6.6cm}@{}}
\caption{Tools and services used to build and publish the repositories.}\label{tab:tools}\\
\toprule Tool & Where & Role and why\\\midrule\endfirsthead
\toprule Tool & Where & Role and why\\\midrule\endhead
QGIS 4 (the 4.0 release named Norrk\"oping in the project files) & the vineyard and slurry repositories & the desktop GIS in which the maps were built; the project and style and metadata files are in the repository.\\
Google Earth Engine & the course & cloud processing of satellite archives.\\
Sentinel-1 and Sentinel-2 imagery & the course and STADOS and the vineyard repositories & the satellite data analysed.\\
Jupyter and nbconvert & the course and STADOS & notebooks and their HTML export.\\
PyInstaller & three desktop programs & builds one Windows file with no Python needed.\\
GitHub Actions & 8 repositories & a syntax check of all Python files on every push (Python 3.11) in seven, and a test and style and coverage run on Python 3.10 and 3.11 and 3.12 in the dairy model.\\
GitHub Pages and \texttt{.nojekyll} & the website and the graphical abstracts & publishes the pages from the repository. Why: free hosting that needs no server.\\
Zenodo and DOI & the articles & permanent archives of data and code.\\
Citation files (\texttt{CITATION.cff}) & 14 repositories & tells readers how to cite the work.\\
Google Fonts & the website and graphical abstracts & the typefaces (IBM Plex and Space Grotesk and Playfair Display and others).\\
Electron and electron-builder and Node.js & SWAFL & the desktop shell and the portable Windows file (see Supplementary Material 5).\\
pandoc and XeLaTeX and BibTeX & SWAFL & the manual and the articles.\\
Playwright & SWAFL & tests of the real program.\\
Word and Excel formats & manuscripts and reports & files for co-authors and reviewers.\\
\bottomrule\end{longtable}

\section{Installing the Python packages of the account}
The package names that \texttt{pip} uses differ from the import names for four of them: scikit-learn (imported as \texttt{sklearn}), Pillow (\texttt{PIL}), python-docx (\texttt{docx}) and PyTorch (\texttt{torch}, installed as \texttt{torch}).
\begin{verbatim}
python3 -m pip install numpy pandas matplotlib scipy scikit-learn joblib pillow openpyxl
python3 -m pip install scikit-image torch torchvision PyQt6 tensorflow
python3 -m pip install pykrige python-docx xlrd pytest flake8 pyinstaller jupyter
\end{verbatim}
The geospatial packages depend on GDAL and install most reliably from conda-forge, as the requirement files of the repositories advise.
\begin{verbatim}
conda install -c conda-forge geopandas rasterio fiona shapely pyproj rasterstats \
    folium owslib gdal
\end{verbatim}
QGIS is a separate program and is installed from \url{https://qgis.org}. Each repository also has its own \texttt{requirements.txt}, which lists only what that repository needs.

\section{What is not claimed}
Versions are not stated for packages whose requirement files give only a minimum or no version, because no exact version is recorded. Two repositories carry the map files of one QGIS project each and no code. The GIS repositories were made in QGIS, so their method is recorded in the project and style files and the README, not in scripts. The names of the physical models in the gully twin (HEC-HMS and HEC-RAS 2D and RUSLE-LISEM) are those named in its README; their programs are not part of the repository.
\end{document}
